% Источники: exam/materials/преподаватель/2026/Факторный_анализ_и_анализ_главных_компонент.pdf; exam/materials/моделирование.pdf, разделы 38 и 40; exam/materials/прошлогодние ответы/Modelirovanie_1_-3.pdf, вопрос 38.
\section{Факторный анализ и анализ главных компонент в математическом моделировании.}

\subsection*{Факторный анализ}

При исследовании сложных объектов и систем часто невозможно непосредственно измерить величины, определяющие свойства этих объектов (факторы), а иногда неизвестны содержание и смысл факторов. Для измерений доступны только величины в той или иной степени, зависящие от факторов.

При этом, когда влияние неизвестного фактора проявляется в нескольких измеряемых признаках, эти признаки могут обнаруживать тесную связь между собой (коррелированность), поэтому общее число факторов может быть гораздо меньшим, чем число измеряемых переменных, которое обычно выбирается исследователем произвольно.

\textbf{Цель факторного анализа:} компактное всестороннее описание объекта исследования и выявление скрытых переменных факторов, отвечающих за наличие линейных статистических корреляций между наблюдаемыми переменными и, как следствие, сокращение количества переменных.

Обязательные условия факторного анализа:
\begin{itemize}
  \item все признаки должны иметь количественное выражение;
  \item число наблюдений должно значимо больше числа переменных;
  \item выборка должна быть однородна, то есть элементы ее принадлежат одной генеральной совокупности;
  \item исходные переменные должны иметь симметричные распределения;
  \item потенциальные факторы должны коррелировать в большей или меньшей степени.
\end{itemize}

\subsection*{Анализ главных компонент}

Выбор новых признаков, которые являются линейными комбинациями прежних и «вбирают» в себя большую часть общей изменчивости наблюдаемых данных, а поэтому передают большую часть информации, заключенной в первоначальных наблюдениях. Обычно для этого используют метод главных компонент.

Этот метод сводится к выбору новой ортогональной системы координат в пространстве наблюдений.

\textbf{Постановка задачи:} найти подпространства меньшей размерности, в ортогональной проекции на которые разброс данных (дисперсия) максимальна.

Алгоритм: пусть дан центрированный набор (из исходного набора вычли мат. ожидание) векторов $X_i$. Ищется такое ортогональное преобразование в новую систему координат, для которого верны условия:
\begin{enumerate}
  \item выборочная дисперсия данных вдоль 1-ой координаты максимальна;
  \item выборочная дисперсия данных вдоль 2-ой координаты максимальна при условии ортогональности 1-ой координате;
  \item и т.д.
\end{enumerate}

После преобразования — компоненты, а не факторы, пространство другое. На практике доказано, что главные векторы равны собственным векторам ковариационной матрицы
\[
  \frac{1}{n-1}X^T X.
\]

\subsection*{Пример}

Пусть после центрирования исходных наблюдений получена матрица $X_c$. Тогда строят ковариационную матрицу
\[
  C=\frac{1}{n-1}X_c^TX_c.
\]
Главные направления находятся как собственные векторы матрицы $C$. В примере из источника для матрицы наблюдений с двумя одинаково изменяющимися признаками одно собственное значение равно нулю, поэтому разброс вдоль соответствующего ортогонального направления равен нулю; первая главная компонента задаётся собственным вектором, отвечающим ненулевому собственному значению.

\clearpage
